\section{Related Work}
\label{sec:related}

\subsection{Foundation VLA Models}

The powerful understanding and generalization capabilities of VLMs have significantly boosted the development of VLA models. A prevalent and effective approach~\citep{rt2, openvla, octo, gr1} involves directly fine-tuning a pre-trained VLM on specific robotics trajectory datasets, teaching it to output corresponding action tokens based on environmental context and language instructions. RT-2~\citep{rt2} tokenizes actions into the same format as text and co-trains the model on a mixture of vision-language and robotic control datasets, enabling it to directly generate action tokens for motor commands. Following this paradigm, OpenVLA~\citep{openvla} further explores efficient fine-tuning strategies for VLA models. Although these end-to-end methods can fulfill simple commands, this approach risks degrading the VLM's native vision-language understanding and reasoning capabilities, often leading to sub-optimal performance on complex tasks.

To mitigate this issue, a decoupled dual-system architecture~\citep{pi0, gr00t, fstd, thinkact, fis, onetwovla} has been widely explored, where ``System-2" performs high-level tasks and scene understanding and ``System-1" translates relevant information into low-level, executable actions. $\pi_0$~\citep{pi0} appends a diffusion-based action expert to a pre-trained VLM, allowing the model to inherit robust vision-language understanding from web-scale data while achieving precise manipulation control. Similarly, GR00T~\citep{gr00t} adopts a comparable architecture and augments its training data by learning latent action representations from demonstration videos, which enhances its performance and generalization. However, these state-of-the-art models typically assume that user instructions are provided once at the beginning of a task and remain static, thereby overlooking the dynamics of real-world human-robot interaction.

\subsection{Interactive VLA Systems}
Considering the fluid nature of human intent and the inherent limitations of current models, the ability to perceive, respond, and execute concurrently, as well as adapt to new instructions timely is crucial. SayCan~\citep{saycan} combines the high-level semantic understanding of large language models with learned robotic affordances, allowing robots to execute complex, multi-step natural language instructions. VILA~\citep{vila} uses vision-language reasoning for long-horizon task planning and integrates natural visual feedback to perform complex long-term tasks. RT-H~\citep{rth} and YAY Robot~\citep{yay} introduce an intermediate-language-based action hierarchy, where high-level policies generate instructions that guide low-level policy execution, permitting human intervention and adjustment via natural language. RACER~\citep{dai2024racer} achieves task execution and error correction through a supervised-executor dual-model architecture. Recently, Hi-Robot~\citep{hirobot} makes progress in this direction by using a high-level VLM to translate multi-stage, complex instructions into simplified atomic steps for a low-level VLA model to execute. During execution, the high-level VLM can process new instructions and adjust the control policy for the next step after the current atomic step has been completed. Another work, Switch-VLA~\citep{switchvla}, incorporates contact state and behavior modes to seamlessly achieve task switching based on the latest instruction.

Nevertheless, the interactivity of these approaches is still constrained. They must complete their current atomic action or inference cycle before processing a new user directive, and often cannot be interrupted mid-action. This introduces significant latency and limits the system's real-time responsiveness and flexibility. Although Switch-VLA achieves faster response by considering language commands at every action generation step, this design constrains the size of the VLM that can be used, ultimately limiting its capabilities and performance. In this paper, inspired by and based on the interruptible full-duplex voice interaction system VITA~\citep{vita, vita15}, we propose VITA-E, a novel VLA framework that supports large-scale models while enabling fluid interaction, concurrent seeing, hearing, speaking, and acting, and the instant interruption of any ongoing task.
